% Problemas del capítulo 12. Generado por tools/problemas/gen.py; no editar a mano.
\section{Problemas}\label{sec:problemas-12}

Graduados: \(\star\) directo, \(\star\star\) para pensar, \(\star\star\star\) desafío. Las respuestas están en el apéndice~\ref{ap:respuestas}; conviene intentarlos antes de mirarlas.

\begin{problema}[\(\star\)]\label{pb:12-1}
Se mide tres veces la misma distancia: \(25{,}432\), \(25{,}428\) y \(25{,}437\) m. ¿Qué valor anotás? Calculá los residuos.
\end{problema}

\begin{problema}[\(\star\star\)]\label{pb:12-2}
Se conoce la cota de \(A\), \(10{,}000\) m. Se miden los desniveles \(A\to B=+1{,}250\), \(B\to C=+0{,}830\) y \(C\to A=-2{,}074\) m. ¿Cuánto no cierra el circuito? Repartí el error en partes iguales y calculá las cotas de \(B\) y \(C\).
\end{problema}

\begin{problema}[\(\star\star\)]\label{pb:12-3}
Ajustá una recta a \((0,0)\), \((1,1)\), \((2,1)\) y \((3,3)\) de dos maneras: midiendo el error en vertical y midiendo el error perpendicular a la recta. Compará las pendientes.
\end{problema}

\begin{problema}[\(\star\star\)]\label{pb:12-4}
Dos puntos tienen coordenadas locales \((0,0)\) y \((100,0)\), y coordenadas generales \((500,200)\) y \((560,280)\). Encontrá la escala, el giro y la traslación, y pasá a coordenadas generales el punto local \((50,50)\).
\end{problema}

\begin{problema}[\(\star\star\star\)]\label{pb:12-5}
Se midieron cuatro puntos del borde de una rotonda: \((10,0)\), \((0,10)\), \((-10,0)\) y \((0;\ -9{,}8)\) m. Ajustá una circunferencia escribiendo \(x^2+y^2=2ax+2by+c\), que es lineal en \(a\), \(b\) y \(c\), y resolviendo por cuadrados mínimos. Calculá el centro y el radio.
\end{problema}
